\chapter{为什么做：业务起点与价值边界}

\section{先拆开动机、证据和假设}

“为什么做这个项目”不能只回答“电商很热门”，也不能把个人经历直接包装成市场需求。更可靠的起点是把三类信息分开。

\begin{factbox}
\textbf{本人确认的动机：}项目负责人以广州实习和后续求职为主要场景，希望同时展示 AI 产品、数据产品与 AI 应用、Agent 工程能力。由于当前实习位于运营中心，项目选择“运营如何从大量评论中形成可复核判断”作为主问题，而不是选择与个人经历脱节的通用聊天机器人。
\end{factbox}

\begin{inferencebox}
\textbf{有限的方向证据：}私有岗位样本中出现过 AI 产品、AI 应用、Agent、电商、跨境运营和供应链相关职责。这支持“该项目与部分目标岗位有交集”，但不能外推为广州完整招聘市场，也不能证明所有电商企业都有相同需求。私有岗位明细不属于公开手册证据。
\end{inferencebox}

\begin{boundarybox}
\textbf{待验证的产品假设：}产品或类目运营可能希望减少重复阅读评论的时间，并让单品诊断、竞品比较和问题分流返回可复核证据。需求强度、使用频率、现有替代流程、团队协作方式和付费意愿仍需访谈或工作流试用验证。
\end{boundarybox}

上述区分见 \filelink{docs/master-plan.md} 的“为什么做这个项目”和 \filelink{docs/interview-casebook.md} 的 Part 0。

\section{项目服务谁}

\subsection{B 端主用户与协作方}

\begin{factbox}
\textbf{主要用户}是耐用消费品品牌方或跨境电商团队中的产品、类目运营。产品、质量和客服团队是证据的协作消费者，但 MVP 不为四类角色分别建立四套独立产品。
\end{factbox}

只选择一个主用户有两个原因。第一，运营人员更直接面对商品范围、评论筛选和竞品比较；第二，单一主用户能控制 MVP 范围，使数据、工具和界面围绕同一工作流闭环。协作角色仍然可以使用运营整理后的证据，例如质量团队继续核查问题，客服团队更新知识内容。

\subsection{三个业务工作流}

\begin{enumerate}
  \item \textbf{单 SKU 体验诊断：}围绕一个指定商品和体验维度，整理支持、反对、冲突或证据不足的评论信号。
  \item \textbf{商品证据比较：}在两个到五个指定商品之间，按同一体验维度比较各自证据，而不是直接宣布“谁最好”。
  \item \textbf{问题信号分流：}根据冻结证据规则，整理哪些评论问题值得先核查。这里的优先级是调查顺序，不是缺陷严重度或财务损失排序。
\end{enumerate}

\begin{beginnerbox}
Excel 或 BI 擅长结构化汇总，普通聊天模型擅长开放表达。本项目关注的是另一层问题：怎样把自然语言任务绑定到冻结商品范围、评论证据、稳定引用和验证状态。它可以与现有工具协作，不需要宣称替代所有运营工具。
\end{beginnerbox}

\section{明确的 C 端与业务结论边界}

\begin{boundarybox}
本项目不是消费者购物助手，不回答“我应该买哪个”。它也不做实时抓取、价格、库存、销量、需求预测或市场规模分析；不能从评论信号推出真实缺陷率、销售损失、利润、市场份额或因果结论。评论中的“情绪”同样不能直接等同于商业影响。
\end{boundarybox}

当前产品只能称为 \textbf{B 端内部决策支持原型}。若要称为成熟企业 SaaS，还缺少企业商品、客服或工单连接器，多租户隔离，RBAC 权限，团队任务流，监控与 SLA，部署安全以及真实用户效果验证。该边界记录于 \filelink{docs/master-plan.md} 和 \filelink{docs/current-stage-status.md}。

\section{业务价值：三层表达}

\begin{longtable}{>{\raggedright\arraybackslash}p{0.18\textwidth}>{\raggedright\arraybackslash}p{0.34\textwidth}>{\raggedright\arraybackslash}p{0.35\textwidth}}
\toprule
价值层 & 能力说明 & 当前证据强度 \\
\midrule
\endfirsthead
\toprule
价值层 & 能力说明 & 当前证据强度 \\
\midrule
\endhead
效率价值 & 减少重复筛选评论、查找同一体验维度和拼接证据的机械工作。 & 合理产品假设；尚无真人任务对照，不能声称已经节省工时。 \\
决策支持价值 & 为单品诊断、竞品比较、质量跟进和客服知识更新提供可回查的评论证据输入。 & 工程能力已实现；是否改善真实决策仍待目标用户验证。 \\
治理价值 & 固定数据快照、证据口径、引用、轨迹和失败关闭规则，降低编造、串商品和不可复核风险。 & 已有阶段工件、自动化门和审计支持，是当前证据最强的一层。 \\
\bottomrule
\end{longtable}

\begin{advicebox}
后续真人验证应先定义基线任务和记录方法，再观察任务完成时间、首次找到有效证据的时间、引用复核情况、结果采纳或修改、失败原因和重复使用意愿。这些是建议的验证指标，不是当前项目结果。
\end{advicebox}

\section{技术价值：为什么不是“模型读评论后写答案”}

技术价值来自一条分层证据链，而不是单个模型名称。

\begin{enumerate}
  \item \textbf{输入可追踪：}来源文件先经过预检、完整性校验和快照冻结；后续阶段不能随意换数据。
  \item \textbf{口径先冻结：}在开发检索和 Agent 前，先固定允许使用的体验维度、证据强度、极性、标签来源和引用规则。
  \item \textbf{本地代码掌握事实：}SQLite 与确定性工具负责商品解析、证据查询和引用；外部模型不直接读取原始评论，也不能执行任意查询。
  \item \textbf{模型输出必须复核：}Structured Output 约束格式，Verifier 再验证模型填写的值是否被当前证据授权；不一致时修复一次或确定性降级。
  \item \textbf{状态与成本可控制：}LangGraph 只用于一次澄清和可恢复状态，checkpoint 保存可序列化状态，请求预算按真实外部尝试单调计数。
  \item \textbf{结论有界发布：}本地敏感轨迹和公开聚合轨迹分开，公开结果还需值级扫描、源码身份和完成审计。
\end{enumerate}

\begin{inferencebox}
这些能力共同展示了产品定义、数据治理、检索评估、Agent 编排、安全验证和工程交付。它适合同时说明 AI 产品判断与 AI 应用工程能力；但这是对项目能力结构的归纳，不是对个人独立完成程度或生产效果的自动证明。
\end{inferencebox}

\section{价值链与证据链怎样对应}

\begin{tabularx}{0.94\linewidth}{>{\raggedright\arraybackslash}p{0.22\textwidth}X>{\raggedright\arraybackslash}p{0.27\textwidth}}
\toprule
业务问题 & 工程控制 & 防止的错误 \\
\midrule
后续分析是否基于同一输入 & 冻结快照、哈希和 manifest & 文件变化后仍把指标放在一起比较 \\
不同团队是否使用同一口径 & 带版本的证据合同 & 模型或开发者临时创造标签 \\
证据是否串商品 & 精确关联、稳定引用和回库校验 & 把另一商品评论当作当前商品依据 \\
模型是否越权 & 固定工具、严格合同和一次工具限制 & 任意查询、虚构标识或重复消费额度 \\
结果是否可公开 & 双轨迹、闭集 DTO 和递归值扫描 & 路径、内部标识、评论或秘密进入公开结果 \\
阶段是否真的完成 & 冻结案例、阻断门和完成审计 & 把测试全绿或文件存在误当最终证明 \\
\bottomrule
\end{tabularx}

\section{面试时怎样准确介绍}

\begin{interviewbox}
可以这样表达：“我目前在运营中心实习，求职方向希望兼顾 AI 产品和 Agent 应用，因此把主用户限定为品牌方或跨境团队的产品、类目运营，围绕单品诊断、商品比较和问题分流构建评论证据原型。模型只做受约束规划和表达，本地工具取证，Verifier 核对引用。P0--P6 已通过工程门，P7 的技术工作台已完成受控验证，但真人目标运营试用和业务效果还没有完成。”
\end{interviewbox}

\begin{boundarybox}
不要说“系统提升了营收”“已经节省大量工时”“企业已经采用”或“Agent 已上线”。更准确的结论是：系统已经实现从有界运营问题到可复核证据的工程闭环；真实效率、采纳和决策改善仍待验证。
\end{boundarybox}
